\documentclass[a4paper]{article}
\usepackage{amsfonts}
\usepackage{amsmath}
\usepackage{amssymb}
\usepackage{graphicx}     

\begin{document}
  
\noindent \large Auteur: Stan Van Campenhout \\
\noindent \large Studierichting:  Informatica\\
\noindent \large Oefeningenreeks 4A nr. 48.1\\

\medskip

\normalsize

$ \displaystyle\int\dfrac{\operatorname{Bgtan}x}{x^2}dx = -\displaystyle\int\operatorname{Bgtan}xd\left(\dfrac{1}{x}\right)
= -\dfrac{\operatorname{Bgtan}x}{x}+\displaystyle\int\dfrac{1}{x}d\left(\operatorname{Bgtan}x\right) $\\

$ = -\dfrac{\operatorname{Bgtan}x}{x}+\displaystyle\int\dfrac{1}{x\left(x^2+1\right)} dx $\\

$ \rightarrow \dfrac{1}{x\left(x^2+1\right)} = \dfrac{A}{x} + \dfrac{Bx+C}{x^2+1} $\\

$ A = \dfrac{1}{1} = 1 $\\

$ B\left(i\right) + C = \dfrac{1}{i} = -i \Rightarrow B = -1, C = 0 $\\

$ = -\dfrac{\operatorname{Bgtan}x}{x}+\displaystyle\int\dfrac{1}{x}dx + \displaystyle\int\dfrac{-x}{1+x^2}dx $\\

$ = -\dfrac{\operatorname{Bgtan}x}{x}+ \ln|x| + -\dfrac{1}{2}\displaystyle\int\dfrac{1}{1+x^2}d\left(x^2\right) $\\

$ = -\dfrac{\operatorname{Bgtan}x}{x}+\ln|x| + -\dfrac{1}{2}\ln|1+x^2| +c$\\

\begin{thebibliography}{999}
%\bibitem{a} Referentie
\end{thebibliography}
\end{document} 
